\subsection{固有値・固有ベクトル}

固有値と固有ベクトルは、行列が表す一次変換の本質的な性質を示す
指標である。

\subsubsection{定義と求め方}

$n$ 次正方行列 $A$ に対し、零ベクトルでないベクトル $\vect{v}$ と
スカラー $\lambda$ が
\begin{equation}
 A\vect{v}=\lambda\vect{v},\qquad \vect{v}\ne\vect{0}
\end{equation}
を満たすとき、$\lambda$ を $A$ の\textbf{固有値}、
$\vect{v}$ を固有値 $\lambda$ に属する $A$ の
\textbf{固有ベクトル}という。固有ベクトルは、一次変換 $A$ によって
向きが変わらず、$\lambda$ 倍されるベクトルである。

上の式は
\[
 (A-\lambda I)\vect{v}=\vect{0}
\]
と書ける。この同次連立方程式が $\vect{v}\ne\vect{0}$ という
非自明解をもつためには、係数行列 $A-\lambda I$ が正則であっては
ならない。したがって
\begin{equation}
 \boxed{\det(A-\lambda I)=0}
\end{equation}
が必要である。この $\lambda$ の方程式を $A$ の
\textbf{固有方程式}という。固有方程式を解いて固有値を求め、
各固有値について $(A-\lambda I)\vect{v}=\vect{0}$ を解けば、
固有ベクトルが得られる。

\begin{example}[固有値と固有ベクトル]
\[
 A=\begin{pmatrix}2&1\\3&4\end{pmatrix}
\]
の固有値と固有ベクトルを求める。固有方程式は
\begin{align*}
 \det(A-\lambda I)
 &=\begin{vmatrix}2-\lambda&1\\3&4-\lambda\end{vmatrix}\\
 &=\lambda^2-6\lambda+5
  =(\lambda-1)(\lambda-5)=0
\end{align*}
であるから、固有値は $\lambda=1,5$ である。

$\lambda=1$ のとき、
\[
 (A-I)\vect{v}
 =\begin{pmatrix}1&1\\3&3\end{pmatrix}
  \begin{pmatrix}x\\y\end{pmatrix}
 =\begin{pmatrix}0\\0\end{pmatrix}
\]
より $x+y=0$ なので、固有ベクトルは
\[
 \vect{v}=c\begin{pmatrix}1\\-1\end{pmatrix}\qquad(c\ne0)
\]
である。同様に、$\lambda=5$ のとき
\[
 (A-5I)\vect{v}
 =\begin{pmatrix}-3&1\\3&-1\end{pmatrix}
  \begin{pmatrix}x\\y\end{pmatrix}
 =\begin{pmatrix}0\\0\end{pmatrix}
\]
より $y=3x$ なので、固有ベクトルは
\[
 \vect{v}=c\begin{pmatrix}1\\3\end{pmatrix}\qquad(c\ne0)
\]
である。
\end{example}

\paragraph{演習}
\[
 A=\begin{pmatrix}3&1\\1&3\end{pmatrix}
\]
の固有値と固有ベクトルを求めよ。

\subsubsection{固有値・固有ベクトルの性質}

\paragraph{異なる固有値に属する固有ベクトル}
異なる固有値に属する固有ベクトルは線形独立である。
実際、$\lambda_1\ne\lambda_2$ に属する固有ベクトルを
$\vect{v}_1,\vect{v}_2$ とし、
\[
 x_1\vect{v}_1+x_2\vect{v}_2=\vect{0}
\]
とする。両辺に $A$ を作用させた式から、元の式を
$\lambda_1$ 倍した式を引くと、
\[
 x_2(\lambda_2-\lambda_1)\vect{v}_2=\vect{0}
\]
となる。$\lambda_2-\lambda_1\ne0$ かつ $\vect{v}_2\ne\vect{0}$
なので $x_2=0$ であり、元の式から $x_1=0$ も従う。
よって $\vect{v}_1,\vect{v}_2$ は線形独立である。
一般に、互いに異なる固有値に属する固有ベクトルの組は線形独立である。

\paragraph{行列式とトレース}
2次正方行列
\[
 A=\begin{pmatrix}a_{11}&a_{12}\\a_{21}&a_{22}\end{pmatrix}
\]
の固有値を $\lambda_1,\lambda_2$ とする。固有多項式は
\begin{align*}
 \det(A-\lambda I)
 &=\lambda^2-(a_{11}+a_{22})\lambda
   +(a_{11}a_{22}-a_{12}a_{21})\\
 &=(\lambda-\lambda_1)(\lambda-\lambda_2)
\end{align*}
である。係数を比較すると
\begin{equation}
 \boxed{\det A=\lambda_1\lambda_2,\qquad
        \operatorname{tr}A=\lambda_1+\lambda_2}
\end{equation}
を得る。一般の $n$ 次正方行列でも、重複度を込めて
\[
 \det A=\prod_{i=1}^n\lambda_i,\qquad
 \operatorname{tr}A=\sum_{i=1}^n\lambda_i
\]
が成り立つ。

\subsubsection{行列の対角化}

$n$ 次正方行列 $A$ が線形独立な固有ベクトル
$\vect{v}_1,\ldots,\vect{v}_n$ をもち、対応する固有値を
$\lambda_1,\ldots,\lambda_n$ とする。固有ベクトルを列に並べた行列
\[
 P=(\vect{v}_1\ \vect{v}_2\ \cdots\ \vect{v}_n)
\]
は正則であり、
\begin{align*}
 AP
 &=(A\vect{v}_1\ A\vect{v}_2\ \cdots\ A\vect{v}_n)\\
 &=(\lambda_1\vect{v}_1\ \lambda_2\vect{v}_2\
    \cdots\ \lambda_n\vect{v}_n)\\
 &=P\begin{pmatrix}
 \lambda_1&&0\\
 &\ddots&\\
 0&&\lambda_n
 \end{pmatrix}
\end{align*}
となる。したがって
\begin{equation}
 \boxed{D=P^{-1}AP
 =\begin{pmatrix}
 \lambda_1&&0\\
 &\ddots&\\
 0&&\lambda_n
 \end{pmatrix}}
\end{equation}
である。このように、正則行列 $P$ によって $P^{-1}AP$ を対角行列に
することを、$A$ を\textbf{対角化する}という。なお、すべての行列が
対角化できるわけではなく、$n$ 個の線形独立な固有ベクトルをもつことが
対角化可能であるための必要十分条件である。

対角化できる行列では $A=PDP^{-1}$ なので、
\begin{equation}
 \boxed{A^m=PD^mP^{-1}}
\end{equation}
となり、高いべきも容易に計算できる。

\begin{example}[対角化による行列のべき]
\[
 A=\begin{pmatrix}2&1\\3&4\end{pmatrix}
\]
について、前の例で求めた固有ベクトルを用いると
\[
 P=\begin{pmatrix}1&1\\-1&3\end{pmatrix},\qquad
 D=\begin{pmatrix}1&0\\0&5\end{pmatrix},\qquad
 P^{-1}=\frac14\begin{pmatrix}3&-1\\1&1\end{pmatrix}.
\]
したがって、非負整数 $m$ に対して
\begin{align}
 A^m
 &=P\begin{pmatrix}1&0\\0&5^m\end{pmatrix}P^{-1}\notag\\
 &=\frac14
 \begin{pmatrix}
 3+5^m&-1+5^m\\
 -3+3\cdot5^m&1+3\cdot5^m
 \end{pmatrix}.
\end{align}
\end{example}

\subsubsection{対称行列の固有値・固有ベクトル}

実対称行列 $A=A^{\mathsf T}$ の固有値はすべて実数であり、異なる
固有値に属する固有ベクトルは互いに直交する。さらに、固有ベクトルを
適切に選ぶことで正規直交基底を作ることができる。したがって、ある
直交行列 $U$ を用いて
\begin{equation}
 \boxed{U^{\mathsf T}AU=D}
\end{equation}
と対角化できる。これは実対称行列の\textbf{スペクトル定理}である。

\begin{example}[対称行列の直交対角化]
\[
 A=\begin{pmatrix}7&2\\2&4\end{pmatrix}
\]
の固有値は $\lambda_1=8,\lambda_2=3$ であり、対応する単位固有ベクトル
として
\[
 \vect{u}_1=\frac1{\sqrt5}\begin{pmatrix}2\\1\end{pmatrix},
 \qquad
 \vect{u}_2=\frac1{\sqrt5}\begin{pmatrix}-1\\2\end{pmatrix}
\]
を選べる。これらは
\[
 \vect{u}_1\mathbin{\boldsymbol{\cdot}}\vect{u}_2=0,\qquad
 \lVert\vect{u}_1\rVert=\lVert\vect{u}_2\rVert=1
\]
を満たす。よって
\[
 U=(\vect{u}_1\ \vect{u}_2)
 =\frac1{\sqrt5}\begin{pmatrix}2&-1\\1&2\end{pmatrix}
\]
は直交行列であり、$U^{-1}=U^{\mathsf T}$ である。実際、
\[
 U^{\mathsf T}AU=\begin{pmatrix}8&0\\0&3\end{pmatrix}.
\]
\end{example}

\subsubsection{2次曲線への応用}

交差項を含む2次曲線
\begin{equation}
 5x^2+2\sqrt3xy+7y^2=4
\end{equation}
を考える。$\vect{x}=(x,y)^{\mathsf T}$ とおくと、
\[
 \vect{x}^{\mathsf T}A\vect{x}=4,\qquad
 A=\begin{pmatrix}5&\sqrt3\\\sqrt3&7\end{pmatrix}
\]
と表せる。$A$ の固有値は $4,8$ であり、対応する単位固有ベクトルを
\[
 \vect{u}_1=\frac12\begin{pmatrix}\sqrt3\\-1\end{pmatrix},
 \qquad
 \vect{u}_2=\frac12\begin{pmatrix}1\\\sqrt3\end{pmatrix}
\]
と選ぶと、
\[
 U=(\vect{u}_1\ \vect{u}_2)
 =\frac12\begin{pmatrix}\sqrt3&1\\-1&\sqrt3\end{pmatrix},
 \qquad
 U^{\mathsf T}AU=\begin{pmatrix}4&0\\0&8\end{pmatrix}=D.
\]
$U$ は角 $-\pi/6$ の回転行列である。座標変換
\[
 \vect{x}=U\vect{x}',\qquad
 \vect{x}'=\begin{pmatrix}x'\\y'\end{pmatrix}
\]
を行うと、
\begin{align*}
 \vect{x}^{\mathsf T}A\vect{x}
 &=(U\vect{x}')^{\mathsf T}A(U\vect{x}')\\
 &=\vect{x}'^{\mathsf T}U^{\mathsf T}AU\vect{x}'
 =\vect{x}'^{\mathsf T}D\vect{x}'
\end{align*}
なので、曲線の方程式は
\[
 4{x'}^2+8{y'}^2=4,
 \qquad\text{すなわち}\qquad
 \boxed{{x'}^2+\frac{{y'}^2}{(1/\sqrt2)^2}=1}
\]
となる。したがって、もとの曲線は主軸の方向が座標軸から
$-\pi/6$ だけ回転した楕円である。

\begin{center}
\centering
\begin{tikzpicture}[>=Stealth,scale=1.65]
 \def\ang{-30}
 \draw[->] (-1.45,0)--(1.55,0) node[right] {$x$};
 \draw[->] (0,-1.15)--(0,1.25) node[above] {$y$};
 \draw[->,blue] (0,0)--(\ang:1.35) node[below right] {$x'$};
 \draw[->,blue] (0,0)--({90+\ang}:1.1) node[above right] {$y'$};
 \draw[rotate=\ang,thick,red]
   plot[domain=0:360,samples=100,smooth]
   ({cos(\x)},{sin(\x)/sqrt(2)});
 \draw[->] (.43,0) arc[start angle=0,end angle=\ang,radius=.43];
 \node[right] at (.48,-.13) {$-\pi/6$};
\end{tikzpicture}

\smallskip
\refstepcounter{figure}\small 図\thefigure\quad
固有ベクトルの方向を主軸とする楕円
\end{center}
